\section{Contributions of the Thesis}\label{sec:intro_contrib}

This thesis does not propose a new learning algorithm or a new explanation
method. Its contributions lie in the formulation, the evaluation design and
the empirical evidence:

\begin{enumerate}
  \item \textbf{Formulation.} First-call causer attribution is formulated
    as semi-inductive link prediction on an event-centric knowledge graph
    built from BMW quality notification, inspection and rework data, with
    an explicit definition of which information is available at the time of
    the first decision and of how the confirmed causer is derived from the
    supplier acknowledgement and internal assignment.

  \item \textbf{Evaluation design.} Test incidents are stratified into
    novelty levels based on how much of their attribute combination has been
    observed before, and the explanations of all three model families are
    converted into sets of graph edges so that their faithfulness can be
    measured with one shared necessary-and-sufficient fidelity protocol.

  \item \textbf{Empirical comparison on industrial data.} A rule-based
    baseline reproducing the deployed repeat-failure search, knowledge graph
    embeddings and a relational graph convolutional network are compared on
    real production data with respect to accuracy, faithfulness and their
    behavior across novelty levels, including how many incidents could be
    attributed early at a given precision.

  % TODO keep item 4 only if the synthetic dataset survives
  \item \textbf{Controlled validation.} The same protocol is applied to a
    synthetic event-centric graph with planted causer rules, which provides
    ground-truth explanations and tests whether the findings depend on
    properties specific to the BMW data.
\end{enumerate}
